\documentclass[12pt]{article}
\usepackage{amsmath}
\usepackage{amsfonts}
\usepackage{amssymb}
\usepackage{geometry}
\usepackage{physics}
\usepackage{tikz}
\geometry{a4paper, margin=0.7in}


\title{Homework assignment for 02YMECH \\ Chapter: Small Oscillations}
\date{}
\begin{document}
\maketitle

\vspace*{-2cm}

\section*{Solutions}

\begin{enumerate}

% 1
\item \textbf{Time to first maximum displacement (undamped oscillator).}

The motion is
\[
x(t) = A \cos(\omega_0 t + \phi),
\]
where $A$, $\omega_0$, and $\phi$ are constants. The velocity is
\[
v(t) = \dot{x}(t) = -A \omega_0 \sin(\omega_0 t + \phi).
\]

Maximum displacement in simple harmonic motion occurs when $|x(t)| = A$, i.e. when the cosine is $\pm 1$. At those instants, the velocity is zero.

We are told the initial position and velocity:
\[
x(0) = x_0 = A \cos \phi,
\qquad
v(0) = v_0 = -A \omega_0 \sin \phi.
\]

From these,
\[
\cos \phi = \frac{x_0}{A},
\qquad
\sin \phi = - \frac{v_0}{A \omega_0}.
\]

We want the first time $t>0$ such that $x(t)$ reaches an \emph{extreme} (turning point), i.e. when $v(t)=0$ \emph{and} $x$ is a maximum in magnitude.

The condition $v(t)=0$ gives
\[
\sin(\omega_0 t + \phi) = 0
\quad \Longrightarrow \quad
\omega_0 t + \phi = n\pi,
\quad n \in \mathbb{Z}.
\]

Thus the candidate times are
\[
t_n = \frac{n\pi - \phi}{\omega_0}.
\]

At those instants,
\[
x(t_n) = A \cos(n\pi) = A(-1)^n.
\]
So for even $n$ ($n=0,2,4,\dots$), $x = +A$; for odd $n$ ($n=1,3,5,\dots$), $x = -A$. Both $\pm A$ are extrema of equal magnitude.

We want the \emph{first} $t>0$. The smallest $t_n>0$ is obtained by choosing the smallest integer $n$ such that
\[
t_n = \frac{n\pi - \phi}{\omega_0} > 0.
\]

This gives the rule:
\[
t_{\text{first}} = \min_{n \in \mathbb{Z}} \left\{ \frac{n\pi - \phi}{\omega_0} \, \Big| \, \frac{n\pi - \phi}{\omega_0} > 0 \right\}.
\]

In words: find the next zero of the sine after $t=0$. Equivalently, if $\phi \in (-\pi,\pi]$,
\[
t_{\text{first}} =
\begin{cases}
\displaystyle \frac{-\phi}{\omega_0}, & \text{if } \phi < 0, \\[1em]
\displaystyle \frac{\pi - \phi}{\omega_0}, & \text{if } \phi \ge 0~.
\end{cases}
\]

We can also express $\phi$ in terms of $x_0,v_0$ using
\[
\phi = \arctan2\!\left(-\frac{v_0}{A\omega_0}, \frac{x_0}{A}\right),
\]
where $\arctan2(y,x)$ is the correct quadrant angle with $\cos\phi = x_0/A$ and $\sin\phi = -v_0/(A\omega_0)$.

Summary: the first time the oscillator reaches an endpoint ($|x|=A$) after $t=0$ is given by the smallest positive member of $t_n = (n\pi-\phi)/\omega_0$.

Physically: the oscillator reaches maximum displacement whenever its velocity first becomes zero.


% 2
\item \textbf{Average kinetic and potential energy over one cycle.}

For a harmonic oscillator of mass $m$ and angular frequency $\omega_0$, we can write
\[
x(t) = A \cos(\omega_0 t + \phi).
\]
Then
\[
v(t) = \dot{x}(t) = -A \omega_0 \sin(\omega_0 t + \phi).
\]

The kinetic energy is
\[
K(t) = \frac{1}{2} m v^2(t)
= \frac{1}{2} m A^2 \omega_0^2 \sin^2(\omega_0 t + \phi).
\]

The potential energy (spring potential) is
\[
U(t) = \frac{1}{2} k x^2(t)
= \frac{1}{2} k A^2 \cos^2(\omega_0 t + \phi),
\]
with $k = m \omega_0^2$ for a simple harmonic oscillator. Thus
\[
U(t) = \frac{1}{2} m \omega_0^2 A^2 \cos^2(\omega_0 t + \phi).
\]

The total mechanical energy is constant:
\[
E = K(t)+U(t)
= \frac{1}{2} m \omega_0^2 A^2.
\]

We now compute time averages over one full period $T = \frac{2\pi}{\omega_0}$:
\[
\langle K \rangle
= \frac{1}{T} \int_0^T K(t)\, dt,
\qquad
\langle U \rangle
= \frac{1}{T} \int_0^T U(t)\, dt.
\]

We use the facts that over any full period,
\[
\langle \sin^2(\omega_0 t + \phi) \rangle = \frac{1}{2},
\qquad
\langle \cos^2(\omega_0 t + \phi) \rangle = \frac{1}{2}.
\]

Therefore,
\[
\langle K \rangle
= \frac{1}{2} m A^2 \omega_0^2 \cdot \frac{1}{2}
= \frac{1}{4} m A^2 \omega_0^2,
\]
\[
\langle U \rangle
= \frac{1}{2} m A^2 \omega_0^2 \cdot \frac{1}{2}
= \frac{1}{4} m A^2 \omega_0^2.
\]

But $E = \tfrac{1}{2} m \omega_0^2 A^2$, so
\[
\langle K \rangle = \frac{E}{2},
\qquad
\langle U \rangle = \frac{E}{2}.
\]

Conclusion: over one complete cycle, the oscillator spends half its energy in kinetic form (on average) and half in potential form. The averages are equal and each is $E/2$.


% 3
\item \textbf{Pendulum with a horizontally accelerating support.}

We have a pendulum of length $L$ and bob mass $m$. Initially it hangs straight down. Then the support point begins accelerating horizontally with constant acceleration $a$ (say to the right). After transients, the pendulum reaches a new equilibrium at some constant angle $\theta$ from the vertical, and then small oscillations about that equilibrium are possible.

We will work in the non-inertial frame accelerating with the support. In that frame, there is an effective ``pseudo-force'' of magnitude $ma$ acting on the bob opposite the acceleration of the support (i.e. to the left if the support accelerates to the right). Gravity $mg$ acts downward.

\begin{enumerate}
    \item \textit{New equilibrium angle $\theta$.}

    In the accelerating frame, equilibrium means the tension $\vec{T}$ in the string balances the vector sum of gravity and the horizontal pseudo-force. Let $\theta$ be the angle the string makes with the vertical, displaced in the \emph{opposite} direction of the pseudo-force so that the bob is ``behind'' the accelerating support.

    Force balance requires that the string aligns with the resultant of $mg$ downward and $ma$ horizontally. Geometrically,
    \[
    \tan \theta = \frac{\text{horizontal pseudo-force}}{\text{gravity}} = \frac{ma}{mg} = \frac{a}{g}.
    \]

    Thus
    \[
    \theta = \arctan\!\left(\frac{a}{g}\right).
    \]

    \item \textit{Period of small oscillations about this equilibrium.}

    For small oscillations about the new equilibrium, the pendulum behaves like a simple pendulum of length $L$ in a uniform effective gravitational field $\vec{g}_\text{eff}$ equal to the vector sum of the real gravity and the pseudo-gravity from acceleration.

    The magnitude of this effective gravitational field is
    \[
    g_\text{eff} = \sqrt{g^2 + a^2}.
    \]

    The restoring torque for small angular displacements $\delta\theta$ about equilibrium is proportional to $mg_\text{eff} L \, \delta\theta$, so the small-angle equation of motion is that of a simple pendulum with $g$ replaced by $g_\text{eff}$.

    Therefore, the small-oscillation angular frequency is
    \[
    \omega_{\text{small}} = \sqrt{\frac{g_\text{eff}}{L}}
    = \sqrt{\frac{\sqrt{g^2 + a^2}}{L}}.
    \]

    The corresponding period is
    \[
    T = \frac{2\pi}{\omega_{\text{small}}}
    = 2\pi \sqrt{\frac{L}{g_\text{eff}}}
    = 2\pi \sqrt{\frac{L}{\sqrt{g^2 + a^2}}}.
    \]

    Check limiting cases:
    \begin{itemize}
        \item If $a=0$, then $g_\text{eff}=g$ and $T=2\pi \sqrt{L/g}$, which is the usual pendulum period.
        \item If $a$ is large, the effective gravity increases, so the period decreases.
    \end{itemize}
\end{enumerate}


% 4
\item \textbf{Damped harmonic oscillator amplitude decay.}

The position of a damped oscillator is
\[
x(t) = A e^{-\beta t} \cos(\omega_1 t + \phi),
\]
where $A$ is the initial amplitude factor, $\beta>0$ is the damping coefficient (with units of 1/time), and $\omega_1$ is the damped angular frequency.

The amplitude envelope of the oscillation is
\[
A_{\text{env}}(t) = A e^{-\beta t}.
\]

We are asked: how long until the amplitude decreases to half its initial value?

We want $t = t_{1/2}$ such that
\[
A_{\text{env}}(t_{1/2}) = \frac{A}{2}.
\]

That is,
\[
A e^{-\beta t_{1/2}} = \frac{A}{2}
\quad\Longrightarrow\quad
e^{-\beta t_{1/2}} = \frac{1}{2}.
\]

Take the natural logarithm:
\[
-\beta t_{1/2} = \ln\!\left(\frac{1}{2}\right) = -\ln 2,
\]
so
\[
t_{1/2} = \frac{\ln 2}{\beta}.
\]

This $t_{1/2}$ is the amplitude ``half-life'' for the damped oscillator. Note that it depends only on $\beta$, not on $\omega_1$ or $A$.
















%5

\item \textbf{Average total mechanical energy over one cycle. Resonant case.}

We consider a standard driven damped harmonic oscillator of mass $m$, spring constant $k$, and damping coefficient $b$ (linear viscous damping force $-b\dot{x}$), driven by an external force
\[
F_\text{drive}(t) = F_0 \cos(\omega t).
\]
The equation of motion is
\begin{equation}
m \ddot{x} + b \dot{x} + k x = F_0 \cos(\omega t).
\label{eq:eom}
\end{equation}
The natural angular frequency is
\[
\omega_0 = \sqrt{\frac{k}{m}},
\]
and the damping rate is
\[
\gamma = \frac{b}{2m}.
\]

In steady state (after transients die out), the motion is sinusoidal at the driving frequency:
\begin{equation}
x(t) = X \cos(\omega t - \delta),
\label{eq:x_of_t}
\end{equation}
where the steady-state amplitude $X$ and phase $\delta$ are
\begin{align}
X &= \frac{F_0/m}{\sqrt{(\omega_0^2 - \omega^2)^2 + (2\gamma \omega)^2}}, \\
\tan \delta &= \frac{2\gamma \omega}{\omega_0^2 - \omega^2}.
\end{align}

We will work in this steady-state regime.





The instantaneous mechanical energy of the mass-spring system is
\begin{equation}
E(t) = K(t) + U(t)
= \frac{1}{2} m \dot{x}^2(t) + \frac{1}{2} k x^2(t).
\label{eq:Et}
\end{equation}

Using \eqref{eq:x_of_t}:
\[
x(t) = X \cos(\omega t - \delta),
\qquad
\dot{x}(t) = -X \omega \sin(\omega t - \delta).
\]

Then
\begin{align}
U(t) &= \frac{1}{2} k X^2 \cos^2(\omega t - \delta), \\
K(t) &= \frac{1}{2} m \omega^2 X^2 \sin^2(\omega t - \delta).
\end{align}

Average over one period $T = 2\pi/\omega$.
We use
\[
\langle \cos^2(\omega t - \delta) \rangle = \langle \sin^2(\omega t - \delta) \rangle = \frac{1}{2}.
\]

Therefore,
\begin{align}
\langle U \rangle 
&= \frac{1}{2} k X^2 \cdot \frac{1}{2}
= \frac{1}{4} k X^2, \\
\langle K \rangle
&= \frac{1}{2} m \omega^2 X^2 \cdot \frac{1}{2}
= \frac{1}{4} m \omega^2 X^2.
\end{align}

Thus, the \emph{cycle-averaged} total mechanical energy is
\begin{equation}
\langle E \rangle
= \langle K \rangle + \langle U \rangle
= \frac{1}{4} k X^2 + \frac{1}{4} m \omega^2 X^2.
\label{eq:avgE_general}
\end{equation}

Using $k = m \omega_0^2$:
\begin{equation}
\langle E \rangle
= \frac{1}{4} m \omega_0^2 X^2 + \frac{1}{4} m \omega^2 X^2
= \frac{m X^2}{4} (\omega_0^2 + \omega^2).
\label{eq:avgE_simplified}
\end{equation}

\textbf{Now evaluate at resonance.}

At \emph{driving frequency equal to the natural frequency}, i.e. $\omega = \omega_0$, the steady-state amplitude becomes
\[
X_\text{res}
= \frac{F_0/m}{\sqrt{(\omega_0^2 - \omega_0^2)^2 + (2\gamma \omega_0)^2}}
= \frac{F_0/m}{2\gamma \omega_0}
= \frac{F_0}{2 m \gamma \omega_0}.
\]

At $\omega = \omega_0$, \eqref{eq:avgE_simplified} gives
\[
\langle E \rangle_{\omega=\omega_0}
= \frac{m X_\text{res}^2}{4} (\omega_0^2 + \omega_0^2)
= \frac{m X_\text{res}^2}{4} (2\omega_0^2)
= \frac{m \omega_0^2 X_\text{res}^2}{2}.
\]

Insert $X_\text{res}$:
\[
\langle E \rangle_{\omega=\omega_0}
= \frac{m \omega_0^2}{2}
\left(\frac{F_0}{2 m \gamma \omega_0}\right)^2
= \frac{m \omega_0^2}{2}
\frac{F_0^2}{4 m^2 \gamma^2 \omega_0^2}
= \frac{F_0^2}{8 m \gamma^2}.
\]

So, at resonance,
\begin{equation}
\boxed{
\langle E \rangle_{\omega=\omega_0}
= \frac{F_0^2}{8 m \gamma^2}
= \frac{F_0^2}{8 m} \left(\frac{2m}{b}\right)^2
= \frac{F_0^2 m}{2 b^2}.
}
\end{equation}

This shows that the average stored mechanical energy at resonance grows as damping becomes smaller (since $\gamma = b/(2m)$).


% 2
\item \textbf{Power delivered by the driving force, and balance with dissipation.}

The instantaneous power delivered by the driving force is
\[
P_\text{drive}(t) = \frac{dW}{dt}=\frac{F_\text{drive}(t)\,dx }{dt} = F_\text{drive}(t)\dot{x}(t)
= F_0 \cos(\omega t)\, \dot{x}(t).
\]

The instantaneous rate of energy loss due to damping is the power associated with the friction force $F_\text{damp} = -b \dot{x}$:
\[
P_\text{damp}(t) = F_\text{damp} \, \dot{x}(t)
= (-b \dot{x}) \dot{x}
= - b \dot{x}^2(t).
\]
This is always $\le 0$: damping removes energy.

We now show that \emph{over one cycle in steady state}, the average power input from the driver equals the average power dissipated by damping.

First, write $\dot{x}(t)$ explicitly in steady state using \eqref{eq:x_of_t}:
\[
x(t) = X \cos(\omega t - \delta),
\quad
\dot{x}(t) = -X \omega \sin(\omega t - \delta).
\]

Thus,
\begin{align}
P_\text{drive}(t)
&= F_0 \cos(\omega t)\, \big[-X \omega \sin(\omega t - \delta)\big].
\end{align}

We can rewrite $\sin(\omega t - \delta)$ in terms of $\sin(\omega t)$ and $\cos(\omega t)$:
\[
\sin(\omega t - \delta)
= \sin(\omega t)\cos\delta - \cos(\omega t)\sin\delta.
\]

Then
\begin{align}
P_\text{drive}(t)
&= -F_0 X \omega \cos(\omega t)\,
\big[ \sin(\omega t)\cos\delta - \cos(\omega t)\sin\delta \big] \\
&= -F_0 X \omega
\left[
\cos(\omega t)\sin(\omega t)\cos\delta
- \cos^2(\omega t)\sin\delta
\right].
\end{align}

Average over one period. We use:
\[
\langle \cos(\omega t)\sin(\omega t) \rangle = 0,
\qquad
\langle \cos^2(\omega t) \rangle = \frac{1}{2}.
\]

So
\begin{align}
\langle P_\text{drive} \rangle
&= -F_0 X \omega
\left[
0 \cdot \cos\delta
- \frac{1}{2}\sin\delta
\right] \\
&= \frac{1}{2} F_0 X \omega \sin\delta.
\label{eq:Pdrive_avg}
\end{align}

Now compute the average power dissipated:
\[
P_\text{damp}(t) = - b \dot{x}^2(t)
= - b X^2 \omega^2 \sin^2(\omega t - \delta).
\]

Then
\[
\langle P_\text{damp} \rangle
= - b X^2 \omega^2 \, \langle \sin^2(\omega t - \delta) \rangle
= - b X^2 \omega^2 \cdot \frac{1}{2}
= -\frac{1}{2} b X^2 \omega^2.
\label{eq:Pdamp_avg_prelim}
\]

We must show that $\langle P_\text{drive} \rangle + \langle P_\text{damp} \rangle = 0$ in steady state (no net change of energy over a cycle).

From steady-state phasor relations, one can show
\[
\sin\delta = \frac{2\gamma \omega}{\sqrt{(\omega_0^2 - \omega^2)^2 + (2 \gamma \omega)^2}}
= \frac{2\gamma \omega}{(F_0/m)/X}.
\]
Equivalently,
\[
F_0 \sin\delta
= 2 m \gamma \omega X
= b \omega X,
\]
since $b = 2 m \gamma$.

Plugging into \eqref{eq:Pdrive_avg}:
\begin{align}
\langle P_\text{drive} \rangle
&= \frac{1}{2} F_0 X \omega \sin\delta \\
&= \frac{1}{2} (b \omega X) X \omega \\
&= \frac{1}{2} b X^2 \omega^2.
\end{align}

So
\[
\boxed{
\langle P_\text{drive} \rangle
= \frac{1}{2} b X^2 \omega^2,
\qquad
\langle P_\text{damp} \rangle
= -\frac{1}{2} b X^2 \omega^2,
}
\]
and therefore
\[
\langle P_\text{drive} \rangle + \langle P_\text{damp} \rangle = 0.
\]

Interpretation: in steady state, over each cycle, the driver supplies exactly the energy that the damper removes. The oscillator's averaged mechanical energy is constant in time even though energy is continuously flowing in and being dissipated as heat.


% 3
\item \textbf{Energy decay, $Q$ factor, and interpretation.}

Now consider a \emph{free} damped harmonic oscillator (no driving force), with weak damping ($\gamma \ll \omega_0$). The displacement is
\[
x(t) = A_0 e^{-\gamma t} \cos(\omega_1 t + \phi),
\]
where $\omega_1 \approx \omega_0$ for light damping, and $\gamma = b/(2m)$.

The mechanical energy of the oscillator at time $t$ is
\[
E(t) = \frac{1}{2} m \dot{x}^2 + \frac{1}{2} k x^2.
\]

For light damping, the energy decays exponentially as
\begin{equation}
E(t) = E_0 \, e^{-2 \gamma t},
\label{eq:Et_decay}
\end{equation}
where $E_0$ is the initial energy.

The \emph{quality factor} $Q$ of a lightly damped oscillator is defined as
\begin{equation}
Q = \frac{\omega_0}{2\gamma}
= \frac{m \omega_0}{b}.
\label{eq:Qdef}
\end{equation}

Solve \eqref{eq:Qdef} for $\gamma$:
\[
\gamma = \frac{\omega_0}{2Q}.
\]

Insert into \eqref{eq:Et_decay}:
\[
E(t) = E_0 \exp\!\left[-2 \left(\frac{\omega_0}{2Q}\right) t \right]
= E_0 \exp\!\left[- \frac{\omega_0 t}{Q} \right].
\]

Thus,
\begin{equation}
\boxed{
E(t) = E_0 \, \exp\!\left( - \frac{\omega_0 t}{Q} \right).
}
\end{equation}

Interpretation:

\begin{itemize}
    \item The total mechanical energy of a lightly damped oscillator decays exponentially in time.
    \item The decay rate is set by $\omega_0/Q$. Large $Q$ means slow energy loss (weak damping); small $Q$ means fast energy loss (strong damping).
    \item Over a timescale $\Delta t \sim Q/\omega_0$, the energy drops by a factor of $1/e$. This timescale is often called the ``ringdown time''.
\end{itemize}



\item \textbf{Phase--space diagrams}

\begin{itemize}
\item \textbf{Left diagram:} Trajectories form closed orbits around the origin. Motion is bounded in time. The equilibrium at the origin is stable.
\item \textbf{Right diagram:} Trajectories move away from the origin. The equilibrium at the origin is \emph{unstable}.
\end{itemize}
\end{enumerate}


\end{document}
